\documentclass[10pt,a4paper]{article}
\usepackage[margin=2cm]{geometry}
\usepackage{amsmath,amssymb,amsthm,graphicx}
\usepackage{hyperref,cleveref}
\newtheorem{theorem}{Theorem}
\newtheorem{lemma}[theorem]{Lemma}
\newtheorem{proposition}[theorem]{Proposition}
\newtheorem{corollary}[theorem]{Corollary}
\newtheorem{definition}[theorem]{Definition}
\newtheorem{example}[theorem]{Example}
\newtheorem{remark}[theorem]{Remark}

\usepackage{hyperref}
\hypersetup{
  citecolor=blue,
  colorlinks=true,
  pdftitle={Böhm-Taylor-dBang}
}

% Unused fbox wrapper omitted.

\def\eg{{\em e.g.}}
\def\cf{{\em cf.}}

%%%%%%%%%%%%%%%%%%%%%%%%%%%%%%%%%%%%%%%%%%%%%%%%%%%%%%%%%%%%%%

\usepackage{amsmath,amssymb}
\usepackage{amsfonts}
\usepackage[utf8]{inputenc}
\usepackage[T1]{fontenc}
\usepackage{thmtools}
\usepackage{amsthm}

%%%%%%%%%%%%%%%%%%%%%%%%%%%%%%%%%%
%   TIKZ
%%%%%%%%%%%%%%%%%%%%%%%%%%%%%%%%%%
\usepackage{tikz}
\usetikzlibrary{calc}
\usetikzlibrary{arrows.meta}
\usetikzlibrary{shapes, patterns, decorations.pathreplacing, decorations.text}
\usetikzlibrary{trees}
\usetikzlibrary{positioning}
\usetikzlibrary{tikzmark, cd}


\usepackage{longtable}
\usepackage{proof}
\usepackage{bussproofs}
% Unused mathpartir wrapper omitted.
\usepackage{dsfont}
\usepackage{amsfonts}
\usepackage{stmaryrd}
\usepackage{mathtools}
\DeclareMathSymbol{\oc}{\mathord}{operators}{`!} % Exact used symbol definition from original CMLL source.
\usepackage{scalerel}


\usepackage{pdflscape}

\usepackage{graphicx}
\usepackage{color}
\usepackage{enumitem}

\usepackage{longtable}
\usepackage{multirow}

\newif\iflong
\longtrue


%%%%%%%%%%%
% Syntaxe %
%%%%%%%%%%%

\newcommand{\<}{\langle}
\renewcommand{\>}{\rangle}
\newcommand{\ctx}[1]{\ensuremath{\langle #1 \rangle}}
\newcommand{\alt}{~\mid~}
\newcommand{\set}[1]{\ensuremath{\{#1\}}}
\newcommand{\der}[1]{\ensuremath{\mathbf{der}(#1)}}
\renewcommand{\deg}[2]{d_{#1}(#2)}

%%%%%%%%%%%%%%%%%%%%%%%%
% Langages - ensembles %
%%%%%%%%%%%%%%%%%%%%%%%%

\newcommand{\proc}{\ensuremath{\mathtt{Proc}}}

%vars
\newcommand{\fv}{\mathbf{fv}}

%dcbn
\newcommand{\dcbn}{\ensuremath{\mathtt{dCbN}}}
\newcommand{\cbn}{\ensuremath{\mathtt{CbN}}}

\newcommand{\dcbnbot}{\ensuremath{\dcbn_\bot}}
\newcommand{\taylordCBN}{\ensuremath{\Lambda_{\mathcal{T}^{\text{dCbN}}}}}

%dcbv
\newcommand{\dcbv}{\ensuremath{\mathtt{dCbV}}}
\newcommand{\cbv}{\ensuremath{\mathtt{CbV}}}


%
\newcommand{\bang}{\ensuremath{\mathtt{Bang}}}
\newcommand{\dbang}{\ensuremath{\mathtt{dBang}}}
\newcommand{\dbangb}{\ensuremath{\dbang_\bot}}

\newcommand{\dbangres}{\ensuremath{\mathtt{\delta Bang}}}

\newcommand{\dbangv}{\ensuremath{\dbang_V}}
\newcommand{\dbangn}{\ensuremath{\dbang_N}}

%Approxs%
\newcommand{\notlambda}{\ensuremath{N_{\lambda x. M}}}

%forme normale
\newcommand{\nf}[1]{\mathbf{nf}(#1)}

%%%%%%%%%%%%%%
% Sémantique %
%%%%%%%%%%%%%%

\newcommand{\interp}[1]{\ensuremath{\bbracket{#1}}}

%%%%%%%%%
% Math  %
%%%%%%%%%

\newcommand{\natT}{\ensuremath{\mathbb{N}}}

%%%%%%%%%%%%%%%%%%%%%%%%%
% Réductions - flèches  %
%%%%%%%%%%%%%%%%%%%%%%%%%

\newcommand{\To}{\ensuremath{\Rightarrow}}

%cbn
\newcommand{\ton}{\ensuremath{\to_{\text{n}}}}
\newcommand{\tov}{\ensuremath{\to_{\text{v}}}}

% cbv
\newcommand{\todcbv}{\ensuremath{\to_{\text{dCbV}}}}


%dbang
\newcommand{\todb}{\mapsto_{\oc}} %réduction de base (non contextuelle)
\newcommand{\todbs}{\to_{\oc s}} %réduction de surface
\newcommand{\todbf}{\to_{\oc}} %réduction complète

%dbang ressources
\newcommand{\Todbr}{\Rightarrow_{\delta}} %(ressources) !!réduc entre termes et sommes
\newcommand{\Todbrs}{\Rightarrow_{\delta s}} %(ressources) !!réduc entre termes et sommes
\newcommand{\Todbrf}{\Rightarrow_{\delta}} %(ressources) !!réduc entre termes et sommes
\newcommand{\todbr}{\to_{\delta}} %(ressources) !!réduc entre termes, pas sommes!!!
\newcommand{\todbrs}{\to_{\delta s}} %(ressources \nativeemdash{} surface) 
\newcommand{\todbrf}{\to_{\delta}} %(ressources \nativeemdash{} full)
\newcommand{\totodbr}{\rightrightarrows_{\delta}} %(ressources \nativeemdash{} parallèle, entre termes)

\newcommand{\up}{\ensuremath{\uparrow}}
\newcommand{\down}{\ensuremath{\downarrow}}

%extension relation à ensembles
\newcommand{\extend}[1]{\overline{#1}}
\newcommand{\extendB}[1]{\ensuremath{#1}^\copyright}
\newcommand{\extendC}[1]{\ensuremath{#1}^c}

%%%%%%%%%%%%%%%%%
% Substitutions %
%%%%%%%%%%%%%%%%%

\newcommand{\es}[2]{[#1/#2]}
\newcommand{\s}[2]{\{#1/#2\}}

%%%%%%%%%%%%%%%%%
% Traductions   %
%%%%%%%%%%%%%%%%%

\newcommand{\tradv}[1]{\ensuremath{{#1}^v}}
\newcommand{\tradn}[1]{\ensuremath{{#1}^n}}


%%%%%%%%%%%%%%%%%
% Taylor & Bohm %
%%%%%%%%%%%%%%%%%

\newcommand{\app}{\triangleleft_\oc} % δbang -> dbang
\newcommand{\appv}{\triangleleft_v} %δbang -> dcbv
\newcommand{\appn}{\triangleleft_n} %δbang -> dcbn

\newcommand{\mcal}[1]{\ensuremath{\mathcal{#1}}}
\newcommand{\setApprox}{\mcal{A}}
\newcommand{\setapprox}{\setApprox}

\newcommand{\taylor}[1]{\ensuremath{\mathcal{T}(#1)}}
\newcommand{\taylorp}[1]{\ensuremath{\mathcal{T}^+(#1)}}
\newcommand{\taylorn}[1]{\ensuremath{\mathcal{T}^n(#1)}}
\newcommand{\taylorv}[1]{\ensuremath{\mathcal{T}^v(#1)}}
\newcommand{\taylornf}[1]{\ensuremath{\text{TNF}(#1)}}
\newcommand{\bt}{\ensuremath{{\tt{AT}}}}
\newcommand{\btb}{\ensuremath{{\tt{AT}_\oc}}}
\newcommand{\btn}{\ensuremath{{\tt{AT}_n}}}
\newcommand{\btv}{\ensuremath{{\tt{AT}_v}}}
\newcommand{\smaller}{\ensuremath{\sqsubseteq}}


\newcommand{\fnl}[1]{\textcolor{black}{#1}}
\newif\iflong\longtrue

\newcommand{\nativeemdash}{\ifmmode\text{\textemdash}\else\textemdash\fi}
\newcommand{\nativeendash}{\ifmmode\text{\textendash}\else\textendash\fi}
